% 第5章 介质纳米共振器中的多极模式动力学
\chapter{介质纳米共振器中的多极模式动力学}
\label{chap:result2}

\section{引言}
\label{sec:r2-intro}

% 介质纳米共振器作为纳米光子学的核心构件
% Mie共振理论简介
% 偶极与四极模式的同时激发与干涉

\section{介质纳米共振器的设计与制备}
\label{sec:r2-design}

% 硅膜上的狭缝共振器设计
% FIB加工参数
% 结构表征

\begin{figure}[htbp]
    \centering
    % 占位符：用户实验数据图
    \fbox{\parbox{0.8\linewidth}{\centering\vspace{3cm} 图5-1：介质纳米共振器\\（实验数据待插入）\vspace{3cm}}}
    \caption{\label{fig:ch5-resonator} 200\,nm厚硅膜上的狭缝共振器。(a) SEM图像；(b) 几何尺寸标注。}
\end{figure}

\section{偶极与四极模式的时域分辨}
\label{sec:r2-multipolar}

% 实验测量的空间-时间数据
% SVD分析确认两个主导模式
% 偶极和四极模式之间的相位延迟

\section{内部波前的传播速度测量}
\label{sec:r2-wavefront}

% 狭缝内传播波前的直接观测
% 相速度与有效折射率的关系
% 亚光速传播的物理机制

\section{数值模拟验证}
\label{sec:r2-simulation}

% FDTD模拟设置
% 模拟与实验的对比
% 模式分析

\begin{figure}[htbp]
    \centering
    % 占位符：用户实验数据图
    \fbox{\parbox{0.8\linewidth}{\centering\vspace{3cm} 图5-2：多极模式时域分辨\\（实验数据待插入）\vspace{3cm}}}
    \caption{\label{fig:ch5-multipolar} 介质纳米共振器的时域分辨测量。(a) 测量的空间-时间能量变化图；(b-c) 选定时刻的场分布；(d) SVD分析的主导模式。}
\end{figure}

\section{本章小结}
\label{sec:r2-summary}

本章研究了硅介质纳米共振器中偶极和四极Mie共振模式的时间干涉动力学。阿秒电子显微术首次在时域直接分辨了两种模式之间的相位延迟，观测到狭缝内亚光速传播的电磁波前。这些结果为Huygens超表面的设计提供了时域层面的物理认知。
